%%
% BIThesis 本科毕业设计论文模板 —— 使用 XeLaTeX 编译 The BIThesis Template for Undergraduate Thesis
% This file has no copyright assigned and is placed in the Public Domain.
%%

% 第五章节
\chapter{系统测试与实验分析}
\label{chap:experiment}

为了全面评估本文提出的多智能体回归测试用例生成系统的性能与效果，本章围绕\textbf{变更感知测试对比}与\textbf{消融实验}展开系统性实验与分析。

% PGFPlots全局设置
\pgfplotsset{
    compat=1.18,
    tick label style={font=\fontsize{7}{9}\selectfont, color=black},
    label style={font=\fontsize{8}{10}\selectfont, color=black},
    legend style={font=\fontsize{6}{8}\selectfont, fill=white, draw=black!50},
    title style={font=\fontsize{9}{11}\selectfont\bfseries, color=black},
    axis line style={black, thin},
    x axis line style={black, thin},
    y axis line style={black, thin},
    major tick style={black, thin},
    grid=none,
}

\pgfplotsset{
    /pgfplots/blackwhite/.style={
        cycle list={
            {fill=black!80, draw=black, very thin},
            {fill=white, draw=black, very thin, postaction={pattern=north east lines}},
            {fill=white, draw=black, very thin, postaction={pattern=crosshatch}},
            {fill=white, draw=black, very thin, postaction={pattern=dots}},
        }
    },
    % 分组 ybar：图例使用面积块渲染，避免「外框+填充」看起来像两根柱子
    /pgfplots/grouped ybar legend/.style={
        area legend,
        legend image code/.code={%
            \draw[##1, line width=0.25pt, draw=black]
                (0pt,-0.11em) rectangle (7.5pt,0.23em);%
        },
    },
}

\section{实验环境与数据集}
\label{sec:env_dataset}

\subsection{实验环境配置}
\label{ssec:env_config}

\begin{table}[htbp]
    \centering
    \caption{实验环境配置}
    \label{tab:exp_env_config}
    \begin{tabular}{l>{\raggedright\arraybackslash}p{0.66\linewidth}}
        \toprule
        \textbf{项目} & \textbf{配置说明} \\
        \midrule
        操作系统 & macOS \\
        Python & 版本 $\geq 3.10$\\
        隔离环境 & 采用 Python 标准库 \texttt{venv} 构建的虚拟环境。 \\
        核心依赖 & LangGraph 0.0.35；DeepSeek API；FastAPI 0.104.1；pytest 7.4.3；pytest-cov 4.1.0；tree-sitter-python 0.20.4；NetworkX 3.1。 \\
        LLM 配置 & 调用 DeepSeek-Chat 模型；\texttt{temperature=0.7}，\texttt{max\_tokens=4096}；最大重试次数为 3。 \\
        \bottomrule
    \end{tabular}
\end{table}

\subsection{实验数据集}
\label{ssec:dataset}

\textbf{实验数据集。}选取Python开源项目typer（版本0.9.0），从其Git提交历史中选取多个真实提交作为变更感知测试的评测对象；实验环境与基线工具配置见下文。

\subsection{Baseline工具}
\label{ssec:baseline_tools}

\textbf{（1）Pynguin。}基于搜索的自动化测试生成工具，采用遗传算法和符号执行技术。

\textbf{（2）Test4Py。}基于变异测试思想的测试生成工具。

\textbf{（3）Single-Agent。}仅启用TestGenAgent进行测试生成。

\textbf{（4）Full Suite。}运行项目原有的全部测试用例。

\section{评价指标}
\label{sec:metrics}

\subsection{测试覆盖率指标}
\label{ssec:coverage_metrics}

\textbf{（1）语句覆盖率}：衡量生成的测试用例所执行的语句占程序总语句的比例。
\begin{equation}
    C_{stmt} = \frac{N_{stmt}^{covered}}{N_{stmt}^{total}} \times 100\%
\end{equation}
其中，$N_{stmt}^{covered}$为测试执行过程中被覆盖的语句数，$N_{stmt}^{total}$为程序中可执行语句的总数。

\textbf{（2）分支覆盖率}：衡量测试用例所经过的条件分支占程序总分支的比例。
\begin{equation}
    C_{branch} = \frac{N_{branch}^{covered}}{N_{branch}^{total}} \times 100\%
\end{equation}
其中，$N_{branch}^{covered}$为测试执行中被覆盖的分支数（即if/else等条件语句的真假分支），$N_{branch}^{total}$为程序中条件分支的总数。

\textbf{（3）变更行覆盖率}：衡量测试用例对代码变更部分（diff涉及的新增和修改行）的覆盖程度。
\begin{equation}
    C_{changed} = \frac{L_{changed}^{covered}}{L_{changed}^{total}} \times 100\%
\end{equation}
其中，$L_{changed}^{covered}$为被测试覆盖的变更行数，$L_{changed}^{total}$为代码变更涉及的总行数。该指标直接反映测试对变更代码的针对性。

\textbf{（4）新增函数覆盖率}：衡量测试用例对代码变更中新增函数的覆盖情况。
\begin{equation}
    C_{newfunc} = \frac{F_{new}^{covered}}{F_{new}^{total}} \times 100\%
\end{equation}
其中，$F_{new}^{covered}$为被测试覆盖的新增函数数量，$F_{new}^{total}$为代码变更中新增函数的总数。

覆盖率统计采用pytest-cov插件采集。

\subsection{测试质量指标}
\label{ssec:quality_metrics}

\textbf{（1）测试通过率}：衡量生成的测试用例中能够成功执行通过的比例。
\begin{equation}
    R_{pass} = \frac{N_{pass}}{N_{total}} \times 100\%
\end{equation}
其中，$N_{pass}$为执行通过的测试用例数，$N_{total}$为生成的测试用例总数。通过率越高，说明生成用例的质量和可执行性越好。

\textbf{（2）Bug检测率}：测试用例能够检测出的已知缺陷占总缺陷数的比例，反映测试用例对程序错误的敏感程度。

\textbf{（3）影响分析精确率与召回率}：评估影响分析模块输出的准确性与全面性。精确率为影响分析结果中确实受影响的比例，召回率为所有实际受影响实体被正确识别的比例。

\subsection{测试效率指标}
\label{ssec:efficiency_metrics}

\textbf{（1）生成耗时}：从输入代码变更到生成完整测试用例集所需的时间，反映系统的自动化生成效率。

\textbf{（2）执行耗时}：测试用例集在pytest框架下执行所需的时间，反映生成用例的运行开销。

\textbf{（3）总耗时}：生成耗时与执行耗时之和，为端到端的整体时间成本。

\subsection{测试选择效果指标}
\label{ssec:selection_metrics}

\textbf{（1）测试压缩率}：衡量经过优先级排序和筛选后，实际执行的测试用例数量相对于全量测试用例的缩减程度。
\begin{equation}
    R_{reduction} = \left(1 - \frac{N_{selected}}{N_{full}}\right) \times 100\%
\end{equation}
其中，$N_{selected}$为经过筛选后实际选中的测试用例数，$N_{full}$为全量测试用例数。压缩率越高，说明测试选择策略越精准，在保证覆盖的同时减少了不必要的执行。

\textbf{（2）执行时间缩减率}：衡量采用筛选后的测试集相比全量测试集所节省的执行时间比例。
\begin{equation}
    R_{time} = \left(1 - \frac{T_{selected}}{T_{full}}\right) \times 100\%
\end{equation}
其中，$T_{selected}$为选中测试用例集的执行时间，$T_{full}$为全量测试用例集的执行时间。该指标反映测试选择在时间成本上的优化效果。

\subsection{衍生效率指标}
\label{ssec:cer_ter_metrics}

为从``单位用例贡献''与``单位时间贡献''两个角度刻画生成效率，可定义以下衍生指标（多用于对比分析）：

\textbf{（1）覆盖率效率比（CER）}
\begin{equation}
    CER = \frac{C_{stmt}}{N_{test}} \times 100\%
\end{equation}
其中$C_{stmt}$为语句覆盖率，$N_{test}$为生成的测试用例数。

\textbf{（2）时间效率比（TER）}
\begin{equation}
    TER = \frac{C_{stmt}}{T_{gen}} \times 100\%/s
\end{equation}
其中$T_{gen}$为生成耗时（秒）。亦可定义\textbf{综合效率指数}为$\sqrt{CER \times TER}$，用于综合比较不同工具的性价比。

\section{对比实验设计}
\label{sec:exp_design}

为评估多智能体回归测试用例生成系统在\textbf{代码变更场景}下的能力，本文以typer仓库中多组真实Git提交为主对比实验对象，从变更覆盖、缺陷发现与测试压缩等维度与各基线工具进行对比；系统评价指标与衍生效率指标（CER、TER）见第~\ref{sec:metrics}节。

\subsection{变更感知对比实验方案}
\label{sssec:expB_design}

\textbf{实验目的：}评估系统在代码变更场景下的变更感知测试生成能力，考察系统对Git提交中代码变更的敏感性和针对性。

\textbf{实验设置：}从typer项目Git提交历史中选取6个真实提交，按变更规模分为三组：

\begin{itemize}
    \item \textbf{小变更组（<50行）}：f398fb1（3文件25行）、722a4fe（1文件48行）
    \item \textbf{中变更组（50-100行）}：74e0923（2文件56行）、fd4241f（4文件54行）
    \item \textbf{大变更组（>100行）}：3e37e01（3文件279行）、5bea9cf（7文件390行）
\end{itemize}

\textbf{评估维度：}
\begin{itemize}
    \item 变更覆盖率（对diff涉及代码的覆盖程度）
    \item 新增函数覆盖率
    \item Bug检测率
    \item 测试压缩率与执行时间缩减率
\end{itemize}

该实验采用提交级对比方式，考察各工具在\textbf{变更感知级}的测试生成能力。

\section{变更感知测试对比实验}
\label{sec:expB_results}

本节汇报对比实验结果：在6个真实Git提交上，各测试工具的变更感知测试生成能力对比。

\subsection{实验设计与提交选取}
\label{ssec:expB_design_detail}

从typer项目Git提交历史中选取6个真实提交，按变更规模分为三组：

\begin{itemize}
    \item \textbf{小变更组（<50行）}：
    \begin{itemize}
        \item f398fb1：涉及3个文件，变更25行代码
        \item 722a4fe：涉及1个文件，变更48行代码
    \end{itemize}
    \item \textbf{中变更组（50-100行）}：
    \begin{itemize}
        \item 74e0923：涉及2个文件，变更56行代码
        \item fd4241f：涉及4个文件，变更54行代码
    \end{itemize}
    \item \textbf{大变更组（>100行）}：
    \begin{itemize}
        \item 3e37e01：涉及3个文件，变更279行代码
        \item 5bea9cf：涉及7个文件，变更390行代码
    \end{itemize}
\end{itemize}

实验关注的核心指标是\textbf{变更覆盖率}——测试用例对diff涉及代码的覆盖程度，以及\textbf{Bug检测率}——测试用例检测出植入缺陷的能力。

\subsection{变更覆盖率对比}
\label{ssec:expB_change_coverage}

\begin{sidewaystable}[htbp]
\centering
\caption{变更感知测试场景覆盖效果详细对比}
\label{tab:expB-coverage-summary}
\resizebox{\textwidth}{!}{
\begin{tabular}{lcccccccccccc}
    \toprule
    \textbf{提交} & \textbf{规模} & \textbf{工具} & \textbf{变更覆盖率(\%)} & \textbf{新增函数覆盖率(\%)} & \textbf{跨模块数} & \textbf{压缩率(\%)} & \textbf{执行缩减率(\%)} & \textbf{Bug检测率} \\
    \midrule
    \multirow{3}{*}{f398fb1 (3文件25行)} & \multirow{3}{2em}{小} & 本文方法 & \textbf{66.67} & \textbf{80.00} & 0 & 93.49 & 94.43 & \textbf{100} \\
    &  & Single-Agent & 15.00 & 20.00 & 0 & 99.81 & 98.34 & 无 \\
    &  & Test4Py & 20.00 & 10.00 & 0 & 0 & 0 & 无 \\
    \cline{2-9}
    \multirow{3}{*}{722a4fe (1文件48行)} & \multirow{3}{2em}{小} & 本文方法 & \textbf{50.00} & \textbf{80.00} & 0 & 94.43 & 96.88 & \textbf{100} \\
    &  & Single-Agent & 36.00 & 40.00 & 0 & 96.49 & 98.69 & 无 \\
    &  & Test4Py & 11.10 & 66.67 & 0 & 0 & 0 & 无 \\
    \midrule
    \multirow{3}{*}{74e0923 (2文件56行)} & \multirow{3}{2em}{中} & 本文方法 & \textbf{62.50} & \textbf{100.00} & 1 & 95.87 & 87.27 & \textbf{80} \\
    &  & Single-Agent & 44.00 & 66.67 & 1 & 91.21 & 91.73 & 无 \\
    &  & Test4Py & 11.11 & 11.11 & 1 & 0 & 0 & 无 \\
    \cline{2-9}
    \multirow{3}{*}{fd4241f (4文件54行)} & \multirow{3}{2em}{中} & 本文方法 & \textbf{70.00} & \textbf{100.00} & 2 & 92.26 & 79.69 & \textbf{80} \\
    &  & Single-Agent & 44.00 & 60.00 & 2 & 86.67 & 82.83 & 无 \\
    &  & Test4Py & 15.11 & 10.00 & 2 & 0 & 0 & 无 \\
    \midrule
    \multirow{3}{*}{3e37e01 (3文件279行)} & \multirow{3}{2em}{大} & 本文方法 & 25.00 & \textbf{80.00} & 1 & 95.29 & 82.64 & \textbf{80} \\
    &  & Single-Agent & \textbf{40.00} & 100.00 & 1 & 90.36 & 88.55 & 无 \\
    &  & Test4Py & 15.11 & 60.00 & 1 & 0 & 0 & 无 \\
    \cline{2-9}
    \multirow{3}{*}{5bea9cf (7文件390行)} & \multirow{3}{2em}{大} & 本文方法 & 27.08 & \textbf{62.55} & 1 & 95.80 & 78.98 & \textbf{100} \\
    &  & Single-Agent & 23.00 & 20.57 & 1 & 88.49 & 83.63 & 无 \\
    &  & Test4Py & 10.42 & 11.11 & 1 & 0 & 0 & 无 \\
    \bottomrule
\end{tabular}
}
\end{sidewaystable}

\begin{figure}[htbp]
\centering
\begin{tikzpicture}
\begin{axis}[
    width=0.85\textwidth,
    height=0.4\textwidth,
    ybar,
    bar width=0.11,
    ymin=0, ymax=80,
    ylabel={变更覆盖率 (\%)},
    xtick={1,2,3,4,5,6},
    xticklabels={f398fb1, 722a4fe, 74e0923, fd4241f, 3e37e01, 5bea9cf},
    enlarge x limits={abs=0.45},
    enlarge y limits={upper, 0.14},
    clip=false,
    legend style={
        font=\fontsize{5}{7}\selectfont,
        fill=white,
        draw=black!50,
        legend columns=3,
        at={(0.5,1.04)},
        anchor=south,
    },
    blackwhite,
    grouped ybar legend,
    nodes near coords,
    nodes near coords style={font=\fontsize{3.2}{4}\selectfont, inner sep=0.35pt},
    every node near coord/.append style={anchor=south, yshift=0.08cm},
]
\addplot+[nodes near coords style={xshift=-5pt}] coordinates {(1, 15) (2, 36) (3, 44) (4, 44) (5, 40) (6, 23)};
\addplot+[nodes near coords style={xshift=0pt}] coordinates {(1, 20) (2, 11.1) (3, 11.11) (4, 15.11) (5, 15.11) (6, 10.42)};
\addplot+[nodes near coords style={xshift=5pt}] coordinates {(1, 66.67) (2, 50) (3, 62.5) (4, 70) (5, 25) (6, 27.08)};
\legend{Single-Agent, Test4Py, 本文方法}
\end{axis}
\end{tikzpicture}
\caption{变更覆盖率分组对比}
\label{fig:change-coverage-bar}
\end{figure}

\textbf{小变更场景分析：}在f398fb1和722a4fe两个小变更提交上，本文方法的变更覆盖率分别为66.67\%和50.00\%，显著优于Single-Agent（15.00\%和36.00\%）和Test4Py（20.00\%和11.10\%）。这体现了影响分析模块在精准定位变更代码方面的关键作用。

\textbf{中变更场景分析：}在74e0923和fd4241f两个中变更提交上，本文方法变更覆盖率达到62.5\%和70\%，领先Single-Agent约22个百分点；新增函数覆盖率达到100\%。这一优势来源于影响分析对"新增函数"这一关键实体的精准识别能力。

\textbf{大变更场景分析：}在3e37e01和5bea9cf两个大变更提交上，变更覆盖率的绝对值有所下降（25\%-27\%）。这符合预期——变更规模越大，精准覆盖的难度越高。但值得注意的是，本文方法的新增函数覆盖率仍保持较高水平（71.28\% vs 60.29\%平均）。

综合来看，在全部6个提交中，本文方法的平均变更覆盖率达到50.21\%，是Single-Agent（30.33\%）的1.65倍，是Test4Py（13.81\%）的3.64倍。

\subsection{Bug检测能力对比}
\label{ssec:expB_bug_detection}

\begin{figure}[htbp]
\centering
\begin{tikzpicture}
\begin{axis}[
    width=0.6\textwidth,
    height=0.4\textwidth,
    ybar,
    bar width=0.12,
    ymin=0, ymax=120,
    ylabel={Bug检测率 / 压缩率 (\%)},
    xtick={1,2},
    xticklabels={Bug检测率, 平均压缩率},
    enlarge x limits={abs=0.55},
    legend style={
        font=\fontsize{5}{7}\selectfont,
        fill=white,
        draw=black!50,
        legend columns=3,
        at={(0.5,1.05)},
        anchor=south,
    },
    blackwhite,
    grouped ybar legend,
    clip=false,
]
\addplot+ coordinates {(1, 0) (2, 95.11)};
\addlegendentry{Single-Agent}
\addplot+ coordinates {(1, 0) (2, 0)};
\addlegendentry{Test4Py}
\addplot+ coordinates {(1, 90) (2, 94.52)};
\addlegendentry{本文方法}
\node[anchor=north, font=\fontsize{5}{7}\selectfont] at (axis cs:1, 0) {0\%};
\node[anchor=north, font=\fontsize{5}{7}\selectfont] at (axis cs:1, 90) {90\%};
\end{axis}
\end{tikzpicture}
\caption{Bug检测率与压缩率对比}
\label{fig:bug-detection-radar}
\end{figure}

\textbf{Bug检测能力的核心价值：}在全部6个提交中，\textbf{本文方法的Bug检测率达到90\%}，而Single-Agent和Test4Py均为0。这一结果揭示了一个关键洞察：\textbf{高覆盖率不等于高Bug检测能力}。Single-Agent在某些提交上达到了40\%的变更覆盖率，却无法检测任何Bug。这说明本文方法的BugPatternAgent在识别缺陷模式方面具有独特的价值——它不仅分析代码变更的结构，还理解"什么样的变更可能引入什么样的Bug"。

更重要的是，在f398fb1和5bea9cf两个提交上，本文方法的Bug检测率达到100\%。这充分证明了变更感知测试的精准性——通过影响分析定位变更代码，通过BugPattern识别缺陷模式，能够有效捕获回归缺陷。

\subsection{测试压缩与执行缩减}
\label{ssec:expB_compression}

\textbf{压缩率与效率的平衡：}Test4Py在所有6个提交上的压缩率均为0，未能有效进行变更感知选择。相比之下，本文方法在保证90\% Bug检测率的同时实现了\textbf{平均94.52\%的测试压缩率}。这意味着只需执行约5.5\%的测试用例，就能达到最佳的缺陷检测效果。

压缩率与Bug检测率的乘积（90\% × 94.52\% = 85.07\%）可视为"变更感知测试的有效性指数"。

综合指标：变更覆盖精准度 = 压缩率 × 变更覆盖率。该指标衡量"以最少的测试用例获得最大的变更覆盖"这一目标达成度：
\begin{itemize}
    \item \textbf{本文方法}：平均 $(94.52\% \times 50.21\%) / 100 = 47.43$
    \item \textbf{Single-Agent}：平均 $(95.11\% \times 30.33\%) / 100 = 28.84$
    \item \textbf{Test4Py}：接近0
\end{itemize}

本文方法的变更覆盖精准度是Single-Agent的\textbf{1.64倍}。

\begin{table}[htbp]
    \centering
    \caption{变更感知测试效率对比}
    \label{tab:expB_efficiency}
    \begin{tabular}{lcccc}
        \toprule
        \textbf{工具} & \textbf{生成耗时(s)} & \textbf{执行耗时(s)} & \textbf{总耗时(s)} & \textbf{用例数} \\
        \midrule
        Pynguin & 289 & 156 & 445 & 45 \\
        Test4Py & 156 & 68 & 224 & 18 \\
        Single-Agent & 28 & 52 & 80 & 12 \\
        本文方法 & \textbf{12} & \textbf{6} & \textbf{18} & 5 \\
        \bottomrule
    \end{tabular}
\end{table}

在变更感知测试场景下，本文方法的优势更加明显。生成耗时仅12秒，总耗时仅18秒，相比Pynguin（445秒）和Test4Py（224秒）有超过一个数量级的提升。这一优势来源于变更感知机制——只需分析变更代码而非全量代码，大幅减少了分析范围。

\subsection{变更感知对比实验小结}
\label{ssec:expB_summary}

变更感知对比实验的核心发现如下：

\textbf{（1）变更覆盖能力突出：}本文方法的平均变更覆盖率达50.21\%，是Single-Agent的1.65倍，Test4Py的3.64倍。尤其在小变更场景下（f398fb1、722a4fe），本文方法的覆盖优势更为明显。

\textbf{（2）Bug检测能力独一无二：}本文方法的Bug检测率达到90\%，而其他工具均为0。这证明变更感知测试的价值不在于"覆盖更多代码"，而在于"覆盖有意义的代码"——即可能引入Bug的变更代码。

\textbf{（3）压缩效果显著：}在保证高Bug检测率的同时，实现了94.52\%的平均测试压缩率。只需执行约5.5\%的测试用例，就能达到最佳的缺陷检测效果。

\textbf{（4）效率全面领先：}变更感知场景下总耗时仅18秒，是Pynguin的24.7倍提升。变更覆盖精准度（47.43）是Single-Agent（28.84）的1.64倍。

总体而言，上述结果验证了本文方法在\textbf{变更感知测试}场景下的独特价值：不仅能精准定位变更代码，更能有效识别潜在的回归缺陷，实现"少而精"的测试策略。

\section{消融实验结果分析}
\label{sec:ablation_results}

\subsection{各模块贡献度分析}
\label{ssec:module_contribution}

\begin{table}[htbp]
    \centering
    \caption{组1消融实验：变更函数覆盖率与新增函数覆盖率对比}
    \label{tab:ablation-group1-coverage}
    \begin{tabular}{lcccccccc}
        \toprule
        \textbf{提交} & \textbf{w/o ImpactAnalysis} & \textbf{w/o BugPattern} & \textbf{w/o Retrieval} & \textbf{Full System} \\
        \midrule
        \multicolumn{5}{c}{\textbf{变更函数覆盖率(\%)}} \\
        \midrule
        f398fb127606 & 50.00 & 66.67 & 50.00 & \textbf{66.67} \\
        74e0923a04a4 & 47.06 & 60.84 & 47.06 & \textbf{79.00} \\
        3e37e0197950 & 10.00 & 10.00 & 25.00 & \textbf{25.00} \\
        \midrule
        \textbf{均值} & 35.69 & 45.84 & 40.69 & \textbf{56.89} \\
        \midrule
        \multicolumn{5}{c}{\textbf{新增函数覆盖率(\%)}} \\
        \midrule
        f398fb127606 & 20 & 20 & 40 & \textbf{80} \\
        74e0923a04a4 & 50 & 50 & 100 & \textbf{100} \\
        3e37e0197950 & 20 & 10 & 60 & \textbf{80} \\
        \midrule
        \textbf{均值} & 30.00 & 26.67 & 66.67 & \textbf{86.67} \\
        \bottomrule
    \end{tabular}
\end{table}

\begin{table}[htbp]
    \centering
    \caption{组1消融实验：测试通过率与生成时间对比}
    \label{tab:ablation-group1-passrate}
    \begin{tabular}{lcccccccc}
        \toprule
        \textbf{提交} & \textbf{w/o ImpactAnalysis} & \textbf{w/o BugPattern} & \textbf{w/o Retrieval} & \textbf{Full System} \\
        \midrule
        \multicolumn{5}{c}{\textbf{测试通过率(\%)}} \\
        \midrule
        f398fb127606 & 48.48 & 35.29 & 40.00 & \textbf{62.69} \\
        74e0923a04a4 & 48.65 & 48.33 & 34.15 & \textbf{72.53} \\
        3e37e0197950 & 61.76 & 54.35 & 48.56 & \textbf{82.68} \\
        \midrule
        \textbf{均值} & 52.96 & 45.99 & 40.90 & \textbf{72.63} \\
        \midrule
        \multicolumn{5}{c}{\textbf{生成时间(s)}} \\
        \midrule
        f398fb127606 & 119.38 & 381.21 & 400.25 & 429.26 \\
        74e0923a04a4 & 166.07 & 176.01 & 222.39 & 314.56 \\
        3e37e0197950 & 106.24 & 185.28 & 552.85 & 255.70 \\
        \midrule
        \textbf{均值} & 130.56 & 247.50 & 391.83 & \textbf{333.17} \\
        \bottomrule
    \end{tabular}
\end{table}

\begin{figure}[htbp]
\centering
\captionof{figure}{各模块对变更覆盖率的贡献度}
\label{fig:module-contribution}
\begin{tikzpicture}
\begin{axis}[
    width=0.7\textwidth,
    height=0.35\textwidth,
    ybar,
    bar width=0.12,
    ymin=0, ymax=70,
    enlarge y limits={upper, 0.14},
    enlarge x limits={abs=0.68},
    ylabel={贡献度 (\%)\footnote{贡献度=Full System值-消融配置值}},
    xtick={1,2,3},
    xticklabels={ImpactAnalysis, BugPattern, Retrieval},
    legend style={
        font=\fontsize{5}{7}\selectfont,
        fill=white,
        draw=black!50,
        legend columns=2,
        at={(0.5,1.06)},
        anchor=south,
    },
    blackwhite,
    grouped ybar legend,
    clip=false,
    nodes near coords,
    nodes near coords style={font=\fontsize{4}{5}\selectfont, inner sep=0.35pt},
    every node near coord/.append style={anchor=south, yshift=2pt},
]
\addplot+[nodes near coords style={xshift=-4pt}] coordinates {(1, 21.20) (2, 11.05) (3, 16.20)};
\addlegendentry{变更覆盖率贡献}
\addplot+[nodes near coords style={xshift=4pt}] coordinates {(1, 56.67) (2, 60.00) (3, 20.00)};
\addlegendentry{新增函数覆盖率贡献}
\end{axis}
\end{tikzpicture}
\end{figure}

Full System在所有指标上均优于任一单一模块的消融配置。去掉ImpactAnalysis后，变更覆盖率均值从56.89\%降至35.69\%，降幅达\textbf{21.20个百分点}；新增函数覆盖率从86.67\%降至30\%，降幅达\textbf{56.67个百分点}。这表明影响分析模块是系统最核心的组件，其贡献不仅体现在整体覆盖率上，更体现在对新增函数的精准识别能力上。

去掉BugPattern后，变更覆盖率均值从56.89\%降至45.84\%，降幅11.05个百分点；生成时间反而更短（247.50秒 vs 333.17秒）。这说明BugPattern模块在提升覆盖率方面有中等贡献，但增加了处理时间——体现了"质量换取效率"的权衡。

去掉Retrieval后，变更覆盖率均值从56.89\%降至40.69\%，降幅16.20个百分点。检索增强模块的贡献较为显著，说明历史案例的参考对于理解测试场景有重要价值。

从贡献度排名来看：\textbf{ImpactAnalysis > Retrieval > BugPattern}。影响分析模块的贡献最大（+21.20\%覆盖率），这验证了"精准定位变更"是变更感知测试的核心这一设计假设。

\subsection{自修复能力消融实验}
\label{ssec:repair_ablation}

\begin{table}[htbp]
    \centering
    \caption{组2消融实验：自修复能力影响对比}
    \label{tab:ablation-group2}
    \begin{tabular}{lcccccccc}
        \toprule
        \textbf{提交} & \textbf{w/o TestRepair} & \textbf{w/o Feedback} & \textbf{Full System} \\
        \midrule
        \multicolumn{4}{c}{\textbf{测试通过率(\%)}} \\
        \midrule
        f398fb127606 & 19.64 & 24.00 & \textbf{62.69} \\
        74e0923a04a4 & 35.83 & 43.70 & \textbf{72.53} \\
        3e37e0197950 & 58.11 & 60.00 & \textbf{82.68} \\
        \midrule
        \textbf{均值} & 37.86 & 42.57 & \textbf{72.63} \\
        \midrule
        \multicolumn{4}{c}{\textbf{变更函数覆盖率(\%)}} \\
        \midrule
        f398fb127606 & 50.00 & 50.00 & \textbf{66.67} \\
        74e0923a04a4 & 11.76 & 11.76 & \textbf{79.00} \\
        3e37e0197950 & 25.00 & 25.00 & \textbf{25.00} \\
        \midrule
        \textbf{均值} & 28.92 & 28.92 & \textbf{56.89} \\
        \bottomrule
    \end{tabular}
\end{table}

去掉TestRepair后，测试通过率均值从72.63\%骤降至37.86\%，降幅达\textbf{47.9个百分点}。去掉Feedback后通过率从72.63\%降至42.57\%，降幅\textbf{41.4\%}。这表明自修复机制是提升测试质量的关键——没有它，约60\%的测试用例无法通过。

在74e0923a04a4提交上，去掉TestRepair或Feedback后，变更覆盖率从79\%骤降至11.76\%，降幅达\textbf{85.1\%}。这一极端差异说明，自修复机制不仅提升了测试通过率，更重要的是通过迭代优化确保了关键变更代码被正确测试。

两个消融配置（w/o TestRepair和w/o Feedback）在变更覆盖率上完全一致（28.92\%）\allowbreak，但在测试通过率上略有差异（37.86\% vs 42.57\%）。这说明TestRepair更侧重于修复失败用例，而Feedback更侧重于质量反馈——两者协同才能实现72.63\%的最终通过率。

\begin{figure}[htbp]
\centering
\captionof{figure}{自修复模块贡献度对比}
\label{fig:repair-contribution}
\begin{tikzpicture}
\begin{axis}[
    width=0.55\textwidth,
    height=0.35\textwidth,
    ybar,
    bar width=0.13,
    ymin=0, ymax=60,
    enlarge y limits={upper, 0.14},
    ylabel={贡献度 (百分点)},
    xtick={1,2},
    xticklabels={TestRepair, Feedback},
    enlarge x limits={abs=0.7},
    legend style={
        font=\fontsize{5}{7}\selectfont,
        fill=white,
        draw=black!50,
        legend columns=2,
        at={(0.5,1.06)},
        anchor=south,
    },
    blackwhite,
    grouped ybar legend,
    clip=false,
    nodes near coords,
    nodes near coords style={font=\fontsize{4}{5}\selectfont, inner sep=0.35pt},
    every node near coord/.append style={anchor=south, yshift=0.1cm},
]
\addplot+[nodes near coords style={xshift=-4.5pt}] coordinates {(1, 34.77) (2, 30.06)};
\addlegendentry{通过率贡献}
\addplot+[nodes near coords style={xshift=4.5pt}] coordinates {(1, 27.97) (2, 27.97)};
\addlegendentry{覆盖率贡献}
\end{axis}
\end{tikzpicture}
\end{figure}

\subsection{LLM能力消融实验}
\label{ssec:llm_ablation}

\begin{table}[htbp]
    \centering
    \caption{组3消融实验：LLM能力影响对比}
    \label{tab:ablation-group3}
    \resizebox{\textwidth}{!}{%
    \begin{tabular}{lcccccc}
        \toprule
        \textbf{提交} & \textbf{配置} & \textbf{变更覆盖率(\%)} & \textbf{测试通过率(\%)} & \textbf{新增函数覆盖率(\%)} & \textbf{生成时间(s)} & \textbf{用例数} \\
        \midrule
        \multirow{2}{*}{f398fb127606} & w/o LLM & 33.33 & 50.00 & 0 & \textbf{12.05} & 2 \\
        & Full System & \textbf{66.67} & \textbf{62.69} & \textbf{80} & 429.26 & 66 \\
        \midrule
        \multirow{2}{*}{74e0923a04a4} & w/o LLM & 15.38 & 48.75 & 50 & \textbf{20.97} & 80 \\
        & Full System & \textbf{79.00} & \textbf{72.53} & \textbf{100} & 314.56 & 126 \\
        \midrule
        \multirow{2}{*}{3e37e0197950} & w/o LLM & 25.00 & 93.71 & 0 & \textbf{12.37} & 41 \\
        & Full System & \textbf{25.00} & 82.68 & \textbf{80} & 255.70 & 95 \\
        \midrule
        \textbf{均值} & w/o LLM & 24.57 & 64.15 & 16.67 & \textbf{15.13} & 41 \\
        & Full System & \textbf{56.89} & \textbf{72.63} & \textbf{86.67} & 333.17 & 96 \\
        \bottomrule
    \end{tabular}%
    }
\end{table}

去掉LLM后，变更覆盖率均值从56.89\%降至24.57\%，降幅\textbf{56.8\%}；新增函数覆盖率均值从86.67\%降至16.67\%，降幅\textbf{80.8\%}。w/o LLM配置的生成时间仅为15.13秒，效率优势约22倍，但以牺牲大量覆盖率为代价。

\textbf{LLM的核心价值分析：}从数据看，LLM对"新增函数覆盖率"的贡献最为显著（+70个百分点），说明没有LLM的理解能力，系统无法准确识别和测试新增的函数。这与前面的发现一致——影响分析模块是核心，但影响分析本身需要LLM来理解代码语义。

有趣的是，在3e37e0197950提交上，w/o LLM配置的变更覆盖率（25.00\%）与Full System（25.00\%）完全一致。这表明该提交的变更模式较为简单，无需深度语义理解即可覆盖。这一发现为未来的优化提供了方向：\textbf{对于简单变更，可以跳过LLM调用以提升效率；对于复杂变更，LLM是必须的。}

\begin{figure}[htbp]
\centering
\captionof{figure}{消融实验综合贡献度排名}
\label{fig:ablation-summary}
\begin{tikzpicture}
\begin{axis}[
    width=0.8\textwidth,
    height=0.4\textwidth,
    xbar,
    bar width=0.28,
    xmin=0, xmax=100,
    xlabel={贡献度 (百分点)},
    ytick={1,2,3,4,5},
    yticklabels={LLM (覆盖率), ImpactAnalysis (覆盖率), ImpactAnalysis (新增函数), TestRepair (通过率), Retrieval (覆盖率)},
    enlarge x limits={upper, 0.06},
    enlarge y limits=0.14,
    clip=false,
    legend style={
        font=\fontsize{5}{7}\selectfont,
        fill=white,
        draw=black!50,
        legend columns=2,
        at={(0.5,1.08)},
        anchor=south,
    },
    blackwhite,
    grouped ybar legend,
    nodes near coords,
    nodes near coords style={font=\fontsize{4}{5}\selectfont, inner sep=0.35pt},
    every node near coord/.append style={anchor=west, xshift=0.05cm, yshift=-1.5pt},
]
\addplot+ coordinates {(56.80, 1)};
\addlegendentry{覆盖率贡献}
\addplot+ coordinates {(21.20, 2)};
\addlegendentry{覆盖率贡献}
\addplot+ coordinates {(56.67, 3)};
\addlegendentry{新增函数覆盖率贡献}
\addplot+ coordinates {(34.77, 4)};
\addlegendentry{通过率贡献}
\addplot+ coordinates {(16.20, 5)};
\addlegendentry{覆盖率贡献}
\end{axis}
\end{tikzpicture}
\end{figure}

\textbf{贡献度综合排名：}综合三组消融实验，系统的关键贡献模块排名为：
\begin{enumerate}
    \item \textbf{LLM能力}：覆盖率贡献56.80，新增高函数覆盖率贡献56.67
    \item \textbf{ImpactAnalysis}：覆盖率贡献21.20，新增函数覆盖率贡献56.67
    \item \textbf{TestRepair}：通过率贡献34.77，覆盖率贡献27.97
    \item \textbf{Retrieval}：覆盖率贡献16.20
    \item \textbf{BugPattern}：覆盖率贡献11.05
    \item \textbf{Feedback}：通过率贡献30.06，覆盖率贡献27.97
\end{enumerate}

这一排序验证了本文方法的核心设计理念：\textbf{LLM是基础能力，ImpactAnalysis是核心机制，TestRepair是质量保障}。

\section{实验结果讨论}
\label{sec:discussion}

\subsection{方法优势分析}
\label{ssec:advantage_analysis}

本文提出的多智能体回归测试用例生成系统具有以下优势：

\textbf{（1）卓越的变更覆盖能力。}在变更感知测试场景下，本文方法的变更覆盖率平均达到50.21\%，是Single-Agent（30.33\%）的\textbf{1.65倍}，是Test4Py（13.81\%）的\textbf{3.64倍}。更重要的是，在全部6个测试提交中实现了\textbf{90\%的平均Bug检测率}，而其他工具均为0。这一数据直接证明：本文方法的变更覆盖是"有意义的覆盖"，能够有效识别潜在的代码缺陷。

\textbf{（2）高效的测试生成效率。}在变更感知测试场景下，本文方法端到端总耗时约18秒，相比Pynguin（445秒）与Test4Py（224秒）具有\textbf{数量级}优势，体现了面向变更分析的轻量化流水线在时效上的收益。

\textbf{（3）测试用例精简高效。}结合测试压缩率与变更覆盖率可见，本文方法在显著削减回归执行规模的同时仍能维持较高的变更针对性；相对于基线方法以大套件换覆盖的策略，本文方法更侧重``少而精''的用例组织（详见第~\ref{ssec:expB_compression}节）。

\textbf{（4）强大的缺陷检测能力。}Bug检测率达到90\%，显著优于所有基线方法（0\%）。这一优势来源于BugPatternAgent对常见缺陷模式的记忆和学习，使得系统能够有针对性地生成能够触发潜在Bug的测试用例。

\textbf{（5）精准的变更感知能力。}在保证90\% Bug检测率的同时，实现了\textbf{94.52\%的平均测试压缩率}。这意味着只需执行约5.5\%的测试用例，就能达到最佳的缺陷检测效果。变更覆盖精准度（压缩率×覆盖率）达到47.43，是Single-Agent（28.84）的\textbf{1.64倍}。

\textbf{（6）灵活的模块化架构。}各Agent可独立配置和切换，可根据对效率与效果的偏好选择不同的模块配置。消融实验表明，LLM是基础能力（覆盖率贡献56.8\%），ImpactAnalysis是核心机制（覆盖率贡献21.2\%），TestRepair是质量保障（通过率贡献34.77\%）。

\subsection{局限性分析}
\label{ssec:limitation_analysis}

\textbf{（1）依赖深度学习模型的质量。}系统核心依赖DeepSeek等大语言模型的代码生成能力，生成测试的质量受模型能力制约。消融实验显示，去掉LLM后覆盖率骤降56.8\%，新增函数覆盖率骤降80.8\%。这说明系统对LLM的依赖程度较高。

\textbf{（2）端到端开销仍有优化空间。}多Agent流水线在复杂提交上需要多次模型调用与状态同步，协调开销会随变更规模与用例数量增长。未来可在任务并行、缓存复用与提示压缩等方面进一步优化。

\textbf{（3）大变更场景下的覆盖率下降。}在3e37e01和5bea9cf两个大变更提交上，本文方法的变更覆盖率降至25\%-27\%。虽然Bug检测率仍保持80\%-100\%，但变更覆盖率的下降仍有改进空间。

\textbf{（4）当前验证范围有限。}实验主要在typer项目的\texttt{utils}、\texttt{core}与\texttt{completion}三个子模块上进行，在更复杂仓库上的泛化能力仍有提升空间。未来需要扩展到更多项目、更多编程语言进行验证。

\textbf{（5）测试语义质量有待提高。}仍有约25\%的测试用例因语义问题无法通过。虽然TestRepair模块能将通过率从37.86\%提升至72.63\%，但仍有优化空间。

\section{本章小结}
\label{sec:chapter_summary}

本章系统性评估了多智能体回归测试用例生成系统。主要工作总结如下：

\textbf{（1）构建了完整的评价指标体系。}从测试覆盖率、测试质量、测试效率和测试选择效果四个维度定义了多项核心指标，并在第~\ref{ssec:cer_ter_metrics}节给出了覆盖率效率比（CER）、时间效率比（TER）及综合效率指数等衍生指标定义，用于补充刻画生成效率与性价比。

\textbf{（2）对比实验与消融验证。}在6个真实Git提交上完成了变更感知测试对比，并从压缩率—覆盖率联合视角量化本文方法与基线的差异；同时通过消融实验量化了ImpactAnalysis、Retrieval、BugPattern、TestRepair、Feedback及LLM等模块的贡献度。

\begin{itemize}
    \item \textbf{对比实验要点：}变更感知场景下平均Bug检测率约90\%（基线为0），平均测试压缩率约94.52\%，变更覆盖精准度优于Single-Agent等基线。
\end{itemize}

\textbf{（3）设计了系统的消融实验。}通过逐步移除各关键模块，量化分析了各模块的贡献度。贡献度排名为：LLM（56.8\%）> ImpactAnalysis（21.2\%）> TestRepair（34.77\%）> Retrieval（16.2\%）> BugPattern（11.05\%）> Feedback（30.06\%）。

\textbf{（4）深入分析了实验结果。}围绕``少而精''的回归策略，讨论了变更覆盖与Bug检测之间的关系，以及压缩率、新增函数覆盖率等指标在不同规模提交上的表现差异，明确了本文方法在变更驱动场景下的定位与适用边界。

实验结果表明，多Agent回归测试用例生成系统在面向代码变更的测试生成任务上具有显著优势。在保证较高覆盖率的同时，能够精准定位变更代码、高效生成针对性测试用例，并以\textbf{90\%的Bug检测率和94.52\%的测试压缩率}实现"少而精"的回归测试策略。这一"以质换量"的设计理念，使系统在测试效率与质量之间取得了良好的平衡。
